\section{Policy Selection and Guarded Rebalancing}
The framework places a policy selector above three concrete runtime mechanisms:
pod migration, client-side connection scaling, and subscription redistribution
across active connections. These mechanisms solve different parts of the same
load-distribution problem. Migration changes where a connection terminates,
scaling changes how much parallelism the client maintains, and rebalancing
changes how subscriptions use that parallelism. Keeping the mechanisms
separate makes each action deployable, while the selected policy decides which
action is appropriate at a given control epoch.

The separation between reward and feasibility is deliberate. Safe-RL research
distinguishes performance optimization from constraint satisfaction
\cite{Gu2024SafeRL}; a migration penalty alone cannot enforce ownership or
readiness. Krasowski et al. distinguish action masking, replacement, and
projection as intervention mechanisms \cite{Krasowski2023ProvablySafe}.
Our masks and execution-time checks restrict the deployable actions, but do
not confer the formal safety guarantees of a verified environment model.

\subsection{Greedy and Learned Policy Roles}
Both policy families use the client-visible measurements and deterministic
execution checks defined above. In the scale-out study, the guarded-greedy
selector requests $a^{\mathrm{mig}}$ whenever that action is admissible and
otherwise stays. Its guard assesses projected source and target pressure
after moving the selected subscription group, so a lightly subscribed but
heavily backlogged candidate need not be accepted. The source-selection and
pressure checks are shared with DQL; they are not learned from DQL's outcomes.

The mixed-workload study uses a different deterministic comparator, labelled
\emph{Heuristic}: it applies threshold rules for subscription-count and backlog
imbalance as well as migration. The independent \emph{Scripted} timing controls
use pressure-unaware schedules and are sensitivity controls, not substitutes
for the stronger guarded-greedy comparison. Table~\ref{tab:policy-roles}
keeps these identities distinct.

\begin{table}[!htbp]
\centering\small
\caption{Policy roles and the comparison each supports. Greedy and DQL share
the rebalancing mechanism; the studies retain their original policy labels.}
\Description{Five policy rows distinguish pressure-aware greedy migration,
the multidimensional heuristic, DQL, random scripted schedules, and unchanged placement.}
\label{tab:policy-roles}
\begin{tabular}{p{0.19\textwidth}p{0.39\textwidth}p{0.32\textwidth}}
\toprule
Policy & Decision rule & Evidence role\\
\midrule
Guarded-greedy & Promote an admissible pressure-improving probe & Strong non-learning scale-out control\\
Heuristic & Threshold rules over count, backlog, and placement & Non-learning multi-objective comparator\\
Frozen DQL & Masked maximum learned action value; no evaluation updates & Learned multi-objective selector\\
Scripted schedules & Fixed or jittered timing; random source and target & Pressure-unaware sensitivity controls\\
Static / scale-out only & Retain established connection assignment & No-rebalancing reference\\
\bottomrule
\end{tabular}
\end{table}

The shared contribution is probe-guided actuation, not the invention of either
greedy search or Q-learning. The experiments compare policy choices without
feeding DQL's holdout decisions back into the scripted rules. Learning traces
can inform separate development work, but automatic DQL-to-script tuning is
not part of the evaluated system.
\subsection{DQL State, Action, and Value Function}

The client-side controller is modeled as a partially observable sequential
decision problem. At each control epoch, client $u$ constructs an observation
vector
\begin{equation}
	o_u(t)=
	\left[
	o^B_u(t),o^N_u(t),o^G_u(t),o^C_u(t),o^P_u(t),o^D_u(t)
	\right],
	\label{eq:dql_observation}
\end{equation}
where $o^B_u(t)$ contains normalized local backlog and backlog growth.
The branch $o^N_u(t)$ describes subscription-count imbalance using terms such as
$N_{u,c}(t)$ and $V^N_u(t)$. The branch $o^G_u(t)$ describes backlog imbalance
using terms such as $G_{u,c}(t)$ and $V^G_u(t)$.
The branch $o^C_u(t)$ contains connection count and effective throughput.
The branch $o^P_u(t)$ contains pressure, subscription count, and cumulative
backlog for the current and probe pods. Finally, $o^D_u(t)$ contains cooldown
and statistics-staleness features. For current-pod features, $c^\star$ denotes
the connection being evaluated. All quantities are normalized before they
enter the Q-network.

The DQL action space is defined using bounded action templates rather than all
possible subscription assignments, which keeps the decision problem tractable:
\begin{equation}
	\mathcal{A}=
	\{a^{\mathrm{stay}},a^{\mathrm{reb}N},a^{\mathrm{reb}G},a^{\mathrm{scale}+},a^{\mathrm{scale}-},a^{\mathrm{mig}}\}.
	\label{eq:dql_actions}
\end{equation}
Here $a^{\mathrm{stay}}$ makes no change. The action $a^{\mathrm{reb}N}$ moves one subscription
from the connection with the most subscriptions to the one with the fewest.
The action $a^{\mathrm{reb}G}$ moves the subscription with the largest predicted decrease
in backlog variance. The action $a^{\mathrm{scale}+}$ opens one new connection and,
when allowed, makes bounded subscription moves toward the target $\Theta_N$.
The action $a^{\mathrm{scale}-}$ closes a subscription-free connection only if the
remaining connections meet the required capacity and observed backlog is zero.
The action $a^{\mathrm{mig}}$ requests promotion of the retained probe. It transfers the
evaluated connection or bounded subscription group only if the admissibility
and hysteresis conditions pass. The DQL network selects
the action template; it does not choose an arbitrary Pod or directly change a
subscription assignment. The deterministic executor computes the concrete
source connection, target connection, and subscription using the load-aware
rules in Algorithms~\suppref{alg:adaptive_parallel_scaling} and
\suppref{alg:connection_rebalancing}, and may reject an inapplicable template as a
guarded no-op.
Here $\Theta_N$ is the configured maximum desired subscription count per
connection, not the network parameter vector $\theta$.
The evaluated legacy checkpoint uses a compatible action profile; both
scaling actions are disabled in the loaded cluster trial. The later
cap-restoring scale-up contract is distinct. Table~\suppref{tab:action-contract-map}
retains these version boundaries and their evaluation scope.
DQL uses two action-value networks with the same architecture but different
roles. The online network $Q_\theta(o,a)$ is the trainable network used for
current action evaluation and epsilon-greedy action selection. The target
network $Q_{\theta^-}(o,a)$ is a delayed copy of the online network. It is not
used for selecting the live action and is not updated by every gradient step;
it is used only to compute a stable Bellman target for the next observation.
For a transition
$\xi_t=(o_u(t),a_u(t),R_u(t),o_u(t+1))$ sampled from replay buffer
$\mathcal{D}$, the target is
\begin{equation}
	y_t =
	R_u(t) + \gamma
	\max_{a'\in\mathcal{A}^{\mathrm{adm}}_u(t+1)}Q_{\theta^-}(o_u(t+1),a'),
	\label{eq:dql_target}
\end{equation}
where $\gamma\in[0,1)$ is the discount factor and $\theta^-$ denotes the target
network parameters. The set $\mathcal{A}^{\mathrm{adm}}_u(t+1)\subseteq\mathcal{A}$
contains the actions admitted by the next observation's execution guards;
an infeasible action cannot supply the bootstrap value. The equation applies
to non-terminal transitions. At a true terminal state, $y_t=R_u(t)$.
An interrupted observation window is not assigned an artificial terminal
reward. The online network parameters $\theta$ are trained by
minimizing
\begin{equation}
	\mathcal{L}(\theta)=
	\mathbb{E}_{\xi_t\sim\mathcal{D}}
	\left[
	\left(y_t-Q_\theta(o_u(t),a_u(t))\right)^2
	\right].
	\label{eq:dql_loss}
\end{equation}
Every $K_{\mathrm{sync}}$ training updates, the target network is synchronized by copying the
online parameters, $\theta^- \leftarrow \theta$. This separation prevents the
same rapidly changing network from defining both the prediction
$Q_\theta(o_u(t),a_u(t))$ and the target $y_t$. The parameters $\theta$ and
$\theta^-$ contain all network weights and biases; $\Theta_N$ is instead a
subscription-count threshold. The learning rate $\alpha$ controls the update
$\theta\leftarrow\theta-\alpha\nabla_\theta\mathcal{L}(\theta)$ and is
distinct from the imbalance weights $\eta_N,\eta_G$.
This formulation makes rebalancing a sequential control problem: a migration,
new connection, or subscription move can reduce future backlog but may incur
immediate migration or connection cost.

\begin{table}[!htbp]
\centering\small
\caption{Configuration of the v3 DQL policy used in the reported holdout,
scale-out simulator, and October live trial. Later learner variants are
reported separately in Appendix~\suppref{app:action-diagnostics}.}
\Description{A reproducibility table of network dimensions, training settings,
reward weights, exploration, and frozen evaluation settings.}
\label{tab:dql-configuration}
\begin{tabular}{p{0.38\textwidth}p{0.54\textwidth}}
\toprule
Setting & Value\\
\midrule
Network; optimizer & Dense 21--32--6; ReLU hidden, linear output; mini-batch gradient descent\\
Discount; learning rate & $\gamma=0.99$; $\alpha=0.001$\\
Replay; batch; target copy & Uniform replay, capacity 5,000; batch 16; copy every 250 updates\\
Bootstrap; loss & Target-network maximum over admitted actions; squared TD error\\
Training budget & 10,000 transitions per candidate; 9,985 updates after batch warm-up\\
Training exploration & $\epsilon_0=1$, multiplier $\exp(\log(0.05)/10{,}000)$; floor 0.05\\
Reward weights $(\lambda_1,\ldots,\lambda_7)$ & $(0.20,0.15,0.20,0.15,0.10,0.05,0.15)$\\
Pressure weights $(\omega_X,\omega_B,\omega_H,\omega_R)$ & $(0.30,0.25,0.20,0.25)$\\
Imbalance weights & $\eta_N=\eta_G=0.5$\\
Frozen evaluation & $\epsilon=0$; no model or replay updates; checkpoint hash retained\\
\bottomrule
\end{tabular}
\end{table}

\subsection{Shared Control Loop}
Algorithm~\ref{alg:compact-control} summarizes the policy-independent runtime.
The full sizing and rebalancing executors and the DQL training loop are
retained in Section~\suppref{app:runtime-details}.

\begin{algorithm}[!htbp]
\small
\caption{Probe-Guided Selection and Guarded Execution}
\label{alg:compact-control}
\SetKwInOut{Input}{Input}
\Input{Active connections, retained probe, declared policy and action profile}
\While{client is active}{
 Read active/probe telemetry and client queue/ACK counters\;
 Construct observation and freshness, capacity, cooldown and handoff masks\;
 Select an admitted action using the declared greedy or DQL policy\;
 Recheck guards before applying the selected action\;
 \eIf{migration is selected and admitted}{
  Prepare standby bindings on the same probe; await receipts\;
  Quiesce the source; wait for fresh zero-in-flight telemetry\;
  Activate target bindings; detach source bindings; retire the source socket\;
  Retain unselected sessions and replenish the empty probe\;
 }{
  Execute the admitted bounded rebalancing/scaling action, or stay\;
 }
 Record the applied action, observed reward and guard outcome\;
 Update learning state only in a declared training run with a valid transition\;
}
\end{algorithm}
\subsubsection{Exploration, Hysteresis, and Migration}

Let $\mathcal{A}^{\mathrm{adm}}_u(t)\subseteq\mathcal{A}$ contain only the templates
that pass the current freshness, capacity, cooldown, and handoff guards. The
DQL policy uses epsilon-greedy exploration over this masked set:
\begin{equation}
	a_u(t)=
	\begin{cases}
		\text{Uniform}(\mathcal{A}^{\mathrm{adm}}_u(t)), & \text{w.p. } \epsilon,\\
		\arg\max_{a\in\mathcal{A}^{\mathrm{adm}}_u(t)} Q_\theta(o_u(t),a), & \text{w.p. } 1-\epsilon.
	\end{cases}
\end{equation}
The exploration probability decreases during training. Frozen validation and
deployment policies use $\epsilon=0$. This prevents unreviewed exploratory
changes during evaluation. Online exploration would require a separate
safety study. Migration is
executed only when the selected action is $a^{\mathrm{mig}}$ and the retained probe
candidate improves the estimated value by a hysteresis threshold:
\begin{equation}
	Q_\theta(o_u(t),a^{\mathrm{mig}}) >
	Q_\theta(o_u(t),a^{\mathrm{stay}})+\delta.
\end{equation}
This threshold prevents oscillatory migrations caused by small or stale
statistical differences. It does not choose the probe endpoint: the Kubernetes
load balancer already fixed that endpoint when the probe socket was opened.
Passing this value test is necessary but not sufficient. The executor also
checks telemetry freshness, target capacity, source--target separation, cooldown,
and the absence of a conflicting handoff. As a result, most probes
can remain telemetry-only and may expire or be recycled without carrying a
subscription.

The exact network input order, normalization constants, and deterministic
comparator thresholds are retained in Section~\suppref{app:feature-order}. Checkpoint persistence and full execution
details are in Section~\suppref{app:runtime-details}.
